%%%%%%%%%%%%%%%%%%%%%%%%%%%%%%%%%%%%%%%%%%%%%%%%%%%%%%%%%%%%%%%%%%%%%%%%%%%
%%%                         INTRODUCAO                                  %%%
%%%%%%%%%%%%%%%%%%%%%%%%%%%%%%%%%%%%%%%%%%%%%%%%%%%%%%%%%%%%%%%%%%%%%%%%%%%

\chapter{Introdução}
\label{ch:introducao}

\setlength{\parskip}{0.3cm}
%\thispagestyle{headings}

O Exame Nacional do Ensino Médio (ENEM) é uma prova aplicada anualmente no
Brasil que tem como objetivo avaliar o desempenho dos estudantes ao final da
educação básica e é organizado pelo Instituto Nacional de Estudos e Pesquisas
Educacionais Anísio Teixeira (INEP), abrangendo quatro áreas do conhecimento,
Linguagens e Códigos, Ciências Humanas, Ciências da Natureza e Matemática, além
da redação do tipo dissertativo-argumentativa. Com a nota obtida, os estudantes
podem concorrer a vagas em universidades públicas, solicitar bolsas em
instituições privadas, financiar os estudos e buscar oportunidades em
universidades internacionais, o que torna o exame uma das principais portas de
entrada para o ensino superior e um dos momentos mais importantes da trajetória
escolar de muitos jovens no país.

O resultado mais recente do Exame Nacional do Ensino Médio (ENEM) de 2024
demonstra que embora a média geral aumentou de 543 para 546 pontos, a área de
Matemática e suas Tecnologias apresentou redução de 535 para 529 pontos
\cite{inep2025enem}, sugerindo que a disciplina de matemática apresenta maior
dificuldade pelos alunos.

A matriz de referência do Exame Nacional do Ensino Médio (ENEM) exige que o
estudante deve selecionar, organizar, relacionar, interpretar dados e
informações, representados de diferentes formas, para tomar decisões e enfrentar
situações problema \cite{matriz_enem}. É destacado por\cite{piton2018analise}
que os estudantes apresentam dificuldades significativas na resolução de itens
que exigem interpretação de enunciados e aplicação de conceitos matemáticos de
forma contextualizada.

De acordo com \cite{soares2026diferencas}, a dificuldade não é determinada
apenas pelo conteúdo matemático, mas de uma combinação de fatores como a
formulação dos enunciados, o grau de contextualização, a complexidade
linguística, o tipo de representação mobilizada e o número de operações
cognitivas exigidas para resolver a questão.

Essas dificuldades destacadas pelos autores representam um desafio para a
educação básica, uma vez que a Matemática desempenha papel fundamental no
desenvolvimento de competências relacionadas à análise, argumentação e resolução
de problemas. Além disso, o baixo desempenho nessa área pode impactar
diretamente o acesso dos estudantes ao ensino superior, considerando a
importância da nota do ENEM nos processos seletivos de universidades. 

Esse cenário torna-se ainda mais desafiador na sociedade atual, moldada por
inovações tecnológicas que transformam os ambientes de aprendizagem, comunicação
e interação social. Essa realidade impõe à área educação a necessidade de investigar
como abordagens pedagógicas mediadas por ferramentas digitais podem ajudar a
superar as dificuldades no ensino de Matemática, tornando fundamental o
desenvolvimento de habilidades analíticas voltadas à resolução de problemas
complexos.

Nesse contexto, destaca-se o Pensamento Computacional (PC), definido por
\cite{wing2006} como uma habilidade analítica essencial para todos os indivíduos
na sociedade moderna, não sendo um conhecimento restrito aos cientistas da
computação. Essa perspectiva foca em estratégias de decomposição de problemas,
reconhecimento de padrões, abstração e elaboração de algoritmos como ferramentas
para o desenvolvimento do raciocínio lógico e da capacidade de solucionar
desafios em diferentes áreas do conhecimento.

No âmbito internacional, a Argentina se destaca como um dos países em que o
avanço do ensino de computação tem sido fortemente impulsionado por resoluções
governamentais. Conforme \cite{brackmann2016computational}, após um manifesto da
Fundação Sadosky, em 2013, o Conselho Federal de Educação estabeleceu a
introdução do ensino de programação como componente obrigatório ou
extracurricular no sistema educacional, medida que impulsionou a criação de
diversas iniciativas voltadas à disseminação de práticas de ensino de computação
no país. 

Segundo \cite{brackmann2016computational}, no Chile, embora o país seja líder em
infraestrutura e uso de Tecnologias da Informação e Comunicação na América
Latina, os estudantes ainda possuem pouco contato com os conceitos formais da
computação. Para mudar esse cenário, os pesquisadores relatam que foi proposta
a inclusão do pensamento computacional no currículo básico.

No Brasil a versão inicial da Base Nacional Comum Curricular (BNCC) não
estabelece a computação como uma área de conhecimento independente.Em vez disso,
o documento apresenta as tecnologias digitais apenas como um "tema integrador". \cite{brackmann2016computational}

A BNCC atual destaca a necessidade de metodologias de ensino que estimulem os
estudantes a fazerem uso sistemático do raciocínio lógico, articulando
matemática ligada ao dia a dia e ao uso de tecnologia digitais.  Esse enfoque
sugere uma mudança notável na direção do ensino, que começa a valorizar mais a
compreensão do que a simples memorização de conteúdos. de raciocinar de maneira
independente e sistemática em situações problemáticas
complexas.\cite{brasil2017}

A BNCC também orienta o uso de tecnologias digitais como recurso pedagógico,
reconhecendo os jogos digitais como instrumentos válidos para o desenvolvimento
de competências cognitivas e para o engajamento dos estudantes. As 
práticas plugadas (uso direto de dispositivos e softwares) são incentivadas
como estratégias para tornar o processo de aprendizagem mais dinâmico e
significativo, aproximando os conteúdos curriculares da realidade digital
vivenciada pelos estudantes. Nesse sentido, os jogos digitais possibilita a simulação de
situações-problema, o raciocínio lógico e o desenvolvimento do PC.\cite{brasil2017}

Assim, conclui-se que a melhoria do desempenho em Matemática no ENEM depende não
apenas da exposição com os conteúdos, mas também da adoção de práticas
pedagógicas que estimulem a participação ativa dos estudantes.